\documentclass[aspectratio=169]{beamer}

% Tema nuevo
\usetheme{CambridgeUS}
\usecolortheme{default}

\usepackage[utf8]{inputenc}
\usepackage[T1]{fontenc}
\usepackage[spanish]{babel}
\usepackage{amsmath,amssymb}
\usepackage{graphicx}
\usepackage{xcolor}

\setbeamertemplate{footline}[frame number]
\setbeamertemplate{navigation symbols}{}

% Paleta UdeC
\definecolor{udecblue}{RGB}{0,60,150}

% Colores generales
\setbeamercolor{structure}{fg=udecblue}

% Titular general (portada y similares)
\setbeamercolor{titlelike}{fg=white,bg=udecblue}

% Títulos de cada diapositiva
\setbeamercolor{frametitle}{fg=white,bg=udecblue}

% Bloques
\setbeamercolor{block title}{bg=udecblue!10,fg=udecblue}
\setbeamercolor{block body}{bg=white,fg=black}

% Datos
\title[Ayudantía 8 -- preCertamen 2]{Ayudantía 8 -- preCertamen 2 \\ \small Pauta Certamen 2}
\author[Ignacio Garrido U.]{Profesor: Cristian Mardones, Ph.D. \\ Ayudante: Ignacio Garrido U.}
\institute[UdeC]{Facultad de Ingeniería \\ Depto. de Ingeniería Industrial \\ Macroeconomía (580323)}
\date[2025-2]{Semestre 2025-2}

\begin{document}

% Portada
\begin{frame}
  \titlepage
\end{frame}

% Qué hemos visto
\begin{frame}{¿Qué hemos visto?}
  \begin{itemize}
    \item Ayudantía 5: Modelo IS--LM en economía cerrada.
    \item Ayudantía 6: Política fiscal y monetaria, caso clásico y trampa de liquidez.
    \item Ayudantía 7: Modelos de economía abierta y Mundell--Fleming.
    \item Ayudantía 8: Repaso general preCertamen 2 (IS--LM, Mundell--Fleming, OA--DA, estática comparativa).
  \end{itemize}
\end{frame}

% =========================
% CERTAMEN 2, 2024
% =========================

\section{Certamen 2, 2024}

% P1 Enunciado
\begin{frame}{Certamen 2 2024 -- Pregunta 1}{Enunciado}
\small
Analice gráficamente la magnitud del impacto que generaría un incremento en el consumo autónomo sobre el PIB, considerando:
\begin{enumerate}
  \item El modelo IS--LM en el caso clásico.
  \item El modelo Mundell--Fleming con tipo de cambio flexible.
  \item El modelo de OA--DA en el largo plazo.
\end{enumerate}
\medskip
A partir de su análisis señale qué modelo exhibe el mayor impacto en el PIB.
\end{frame}

% P1 Blanco
\begin{frame}{Certamen 2 2024 -- Pregunta 1}
  \vfill
  \centering
  \vfill
\end{frame}

% P1 Pauta
\begin{frame}{Certamen 2 2024 -- Pregunta 1}{Pauta conceptual}
\small
\textbf{1. IS--LM (caso clásico)}
\begin{itemize}
  \item $Y$ determinado por el producto de pleno empleo $Y^{PE}$ (OA de LP vertical).
  \item Aumento de consumo autónomo $\Rightarrow$ IS se desplaza a la derecha.
  \item LM fija $\Rightarrow$ ajuste solo en $i$.
  \[
  \Delta Y = 0, \quad \Delta i > 0.
  \]
\end{itemize}

\textbf{2. Mundell--Fleming, Tipo de cambio flexible}
\begin{itemize}
  \item Movilidad perfecta de capitales $\Rightarrow i = i^*$.
  \item $\uparrow C_0 \Rightarrow$ IS se desplaza a la derecha $\Rightarrow i \uparrow$ en primera instancia.
  \item Entrada de capitales $\Rightarrow$ apreciación ($e \downarrow$) $\Rightarrow XN \downarrow$ $\Rightarrow$ IS vuelve.
  \[
  \Delta Y = 0, \quad \Delta i = 0.
  \]
\end{itemize}
\end{frame}

\begin{frame}{Certamen 2 2024 -- Pregunta 1}{Pauta conceptual (cont.)}
\small
\textbf{3. OA--DA en el largo plazo}
\begin{itemize}
  \item Aumento de $C_0$ desplaza DA a la derecha.
  \item \textbf{Corto plazo:} $Y \uparrow$, $P \uparrow$ con OA creciente.
  \item \textbf{Largo plazo:} OA de LP vertical en $Y^{PE}$; el ajuste final es:
  \[
  \Delta Y_{LP} = 0,\quad \Delta P_{LP} > 0.
  \]
\end{itemize}

\textbf{Conclusión (nunca está demás):}
\begin{itemize}
  \item En LP, en los tres marcos $\Delta Y = 0$.
  \item Considerando el \textbf{corto plazo}, el modelo OA--DA con OA creciente muestra el mayor impacto inicial en el PIB.
\end{itemize}
\end{frame}

% P2 Enunciado
\begin{frame}{Certamen 2 2024 -- Pregunta 2}{Enunciado}
\small
A partir de las siguientes ecuaciones del modelo IS--LM, formule la diferencial total y obtenga, a través de estática comparativa, el efecto sobre el PIB y tasa de interés provocada por un alza en las transferencias ($TR$):
\[
IS:\quad Y = [C + cTR + c(1 - t)Y] + (I - bi) + G
\]
\[
LM:\quad \frac{M}{P} = kY - hi
\]
\end{frame}

% P2 Blanco
\begin{frame}{Certamen 2 2024 -- Pregunta 2}
  \vfill
  \centering
  \vfill
\end{frame}

% P2 Pauta
\begin{frame}{Certamen 2 2024 -- Pregunta 2}{Pauta con diferenciales}
\small
\textbf{1. Reordenar IS}
\[
Y = C + cTR + c(1-t)Y + I - bi + G
\]
\[
Y - c(1-t)Y = C + cTR + I + G - bi
\]
\[
\theta Y = A + cTR - bi,\quad \theta = 1 - c(1-t).
\]

\textbf{2. Diferenciales}
\[
\theta\, dY = c\,dTR - b\,di
\]
\[
0 = k\,dY - h\,di \Rightarrow di = \frac{k}{h}dY.
\]

Sustituyendo:
\[
\theta dY = c\,dTR - b\left(\frac{k}{h}dY\right)
\Rightarrow dY\left(\theta + \frac{bk}{h}\right) = c\,dTR.
\]
\[
\boxed{dY = \dfrac{c}{\theta + \frac{bk}{h}}\,dTR > 0.}
\]
\end{frame}

\begin{frame}{Certamen 2 2024 -- Pregunta 2}{Pauta con diferenciales (cont.)}
\small
De LM:
\[
di = \frac{k}{h}dY = \frac{k}{h}\,\dfrac{c}{\theta + \frac{bk}{h}}\,dTR.
\]
\[
\boxed{di > 0 \text{ ante } dTR > 0.}
\]

\textbf{Interpretación:}
\begin{itemize}
  \item $\uparrow TR \Rightarrow$ mayor ingreso disponible $\Rightarrow$ IS se desplaza a la derecha.
  \item En el nuevo equilibrio: $Y \uparrow$ e $i \uparrow$.
  \item Tamaño del efecto depende de $c$, $t$, $b$ y de la pendiente de LM ($k,h$).
\end{itemize}
\end{frame}

% P3 Enunciado
\begin{frame}{Certamen 2 2024 -- Pregunta 3}{Enunciado}
\small
Suponga una economía cerrada:
\[
C = 200 + 0{,}6Y_D,\quad t = 0{,}2,\quad TR = 50,\quad I = 800 - 1000i,
\]
\[
G = 500,\quad L = 2Y - 8000i,\quad \frac{M^S}{P} = 1600.
\]
\begin{itemize}
  \item[a)] Encuentre $C$, $Y$, $i$, $I$ y el balance fiscal en el escenario base.
  \item[b)] Determine el nuevo equilibrio si $G$ se eleva a $800$.
\end{itemize}
\end{frame}

% P3 Blanco
\begin{frame}{Certamen 2 2024 -- Pregunta 3}
  \vfill
  \centering
  \vfill
\end{frame}

% P3 Pauta base
\begin{frame}{Certamen 2 2024 -- Pregunta 3}{Pauta -- Escenario base $G=500$}
\small
\textbf{1. Ingreso disponible y consumo}
\[
T = 0{,}2Y,\quad Y_D = Y - T + TR = 0{,}8Y + 50.
\]
\[
C = 200 + 0{,}6Y_D = 230 + 0{,}48Y.
\]

\textbf{2. IS:}
\[
Y = C + I + G = (230 + 0{,}48Y) + (800 - 1000i) + 500
\Rightarrow 0{,}52Y = 1530 - 1000i.
\]

\textbf{3. LM:}
\[
1600 = 2Y - 8000i \Rightarrow i = \frac{2Y - 1600}{8000}.
\]

Sustituyendo y resolviendo:
\[
Y \approx 2246{,}8,\quad i \approx 0{,}36.
\]
\end{frame}

\begin{frame}{Certamen 2 2024 -- Pregunta 3}{Pauta -- Escenario base $G=500$ (cont.)}
\small
\textbf{4. Variables de interés}
\[
Y_D \approx 0{,}8(2246{,}8) + 50,\quad
C \approx 1308{,}4,
\]
\[
I = 800 - 1000i \approx 438{,}3,
\]
\[
T = 0{,}2Y \approx 449{,}4,\quad
BF = T - G - TR \approx -100{,}6.
\]

\textbf{Conclusión:}
\begin{itemize}
  \item PIB moderado, tasa de interés positiva.
  \item Déficit fiscal cercano a 100.
\end{itemize}
\end{frame}

% P3 Pauta G=800
\begin{frame}{Certamen 2 2024 -- Pregunta 3}{Pauta -- Escenario $G=800$}
\small
\textbf{1. Nueva IS con $G=800$}
\[
Y = (230 + 0{,}48Y) + (800 - 1000i) + 800
\Rightarrow 0{,}52Y = 1830 - 1000i.
\]

LM se mantiene:
\[
i = \frac{2Y - 1600}{8000}.
\]

Resolviendo:
\[
Y \approx 2636{,}4,\quad i \approx 0{,}46.
\]

\textbf{2. Variables de interés}
\[
C \uparrow,\quad I \downarrow \text{ (crowding out)},\quad BF \text{ más deficitario}.
\]

\textbf{Mensaje:} La expansión fiscal aumenta $Y$ y $i$, reduce inversión privada y deteriora el balance fiscal.
\end{itemize}
\end{frame}

% P4 Enunciado
\begin{frame}{Certamen 2 2024 -- Pregunta 4}{Enunciado}
\small
Tipo de cambio flexible y perfecta movilidad de capitales.

\textbf{Bienes:}
\[
C = 800 + 0{,}75Y_D,\quad T = 0{,}2Y,\quad TR = 600,
\]
\[
I = 600 - 500i,\quad G = 800,
\]
\[
X = 300 + 3e\frac{P^f}{P} + 0{,}01Y^f,\quad
M = 250 - 40e\frac{P^f}{P} + 0{,}15Y,
\]
\[
P^f = 1,\quad Y^f = 100000.
\]

\textbf{Dinero:}
\[
L^d = Y - 15000i,\quad M^s = 3000,\quad P = 1,\quad i^f = 0{,}3.
\]

Calcule $Y,C,X,M,I,BF,e$ y luego el nuevo equilibrio si $M^s = 3500$.
\end{frame}

% P4 Blanco
\begin{frame}{Certamen 2 2024 -- Pregunta 4}
  \vfill
  \centering
  \vfill
\end{frame}

% P4 Pauta inicial
\begin{frame}{Certamen 2 2024 -- Pregunta 4}{Pauta -- Caso $M^s=3000$}
\small
\textbf{1. Paridad de tasas}
\[
i = i^f = 0{,}3.
\]

\textbf{2. Mercado de dinero}
\[
3000 = Y - 15000(0{,}3) \Rightarrow Y = 7500.
\]

\textbf{3. Componentes internos}
\[
T = 0{,}2Y = 1500,\quad
Y_D = 7500 - 1500 + 600 = 6600.
\]
\[
C = 800 + 0{,}75(6600) = 5750,
\quad
I = 600 - 500(0{,}3) = 450.
\]
\end{frame}

\begin{frame}{Certamen 2 2024 -- Pregunta 4}{Pauta -- Caso $M^s=3000$ (cont.)}
\small
\textbf{4. Sector externo}
\[
X = 300 + 3e + 0{,}01(100000) = 1300 + 3e,
\]
\[
M = 250 - 40e + 0{,}15(7500) = 1375 - 40e,
\]
\[
XN = X - M = -75 + 43e.
\]

\textbf{5. Equilibrio de bienes}
\[
Y = C + I + G + XN
\Rightarrow 7500 = 5750 + 450 + 800 + (-75 + 43e).
\]
\[
43e = 575 \Rightarrow e \approx 13{,}37.
\]

\textbf{A partir de ahí:} se calculan $X,M,XN$ y $BF$ numéricamente como en la pauta escrita.
\end{frame}

% P4 Pauta Ms=3500
\begin{frame}{Certamen 2 2024 -- Pregunta 4}{Pauta -- Caso $M^s=3500$}
\small
\textbf{1. Nuevo equilibrio monetario}
\[
3500 = Y - 15000(0{,}3) \Rightarrow Y = 8000.
\]

\textbf{2. Ajustes internos}
\[
T = 0{,}2Y = 1600,\quad
Y_D = 8000 - 1600 + 600 = 7000,
\]
\[
C = 800 + 0{,}75(7000) = 6050,\quad
I = 600 - 500(0{,}3) = 450.
\]

\textbf{3. Sector externo (esquema)}
\[
X = 1300 + 3e,\quad
M = 250 - 40e + 0{,}15(8000),
\]
\[
XN = X - M,\quad
Y = C + I + G + XN.
\]

\textbf{Mensaje clave:} $\uparrow M^s$ con TC flexible:
\begin{itemize}
  \item No cambia $i$ (sigue en $i^*$).
  \item Aumenta $Y$ y provoca depreciación ($e \uparrow$).
  \item Mejora $XN$ y la recaudación fiscal.
\end{itemize}
\end{frame}

% =========================
% CERTAMEN 2, 2023
% =========================

\section{Certamen 2, 2023}

% 2023 P1 Enunciado
\begin{frame}{Certamen 2 2023 -- Pregunta 1}{Enunciado}
\small
Suponga una economía:
\[
C = 600 + 0{,}7(Y - T + TR),\quad
I = 350 - 15i,\quad
G = 300,\quad
XN = 100,
\]
\[
TR = 100,\quad
T = 0{,}2Y,\quad
L^D = 0{,}5Y - 10i,\quad
\frac{M}{P} = 1500.
\]

\begin{enumerate}[label=\alph*)]
  \item Encuentre el nivel de producción final y la tasa de interés en el equilibrio de corto plazo.
  \item ¿Cómo cambian estos valores si la inversión autónoma baja de 350 a 300?
\end{enumerate}
\end{frame}

% 2023 P1 Blanco
\begin{frame}{Certamen 2 2023 -- Pregunta 1}
  \vfill
  \centering
  \vfill
\end{frame}

% 2023 P1 Pauta
\begin{frame}{Certamen 2 2023 -- Pregunta 1}{Pauta resumida}
\small
\textbf{(a) Equilibrio inicial}
\[
Y_D = 0{,}8Y + 100,\quad
C = 670 + 0{,}56Y.
\]
\[
Y = C + I + G + XN
\Rightarrow
Y = 1420 + 0{,}56Y - 15i
\Rightarrow 0{,}44Y = 1420 - 15i.
\]
LM:
\[
1500 = 0{,}5Y - 10i.
\]

Resolviendo el sistema:
\[
Y \approx 3084,\quad i \approx 4{,}20.
\]

\textbf{(b) Menor inversión autónoma $I = 300 - 15i$}
\[
0{,}44Y = 1370 - 15i,
\quad
1500 = 0{,}5Y - 10i.
\]
\[
Y \approx 3042,\quad i \approx 2{,}10.
\]

\textbf{Comentario:} pequeña caída de $Y$, fuerte caída de $i$.
\end{frame}

% 2023 P2 Enunciado
\begin{frame}{Certamen 2 2023 -- Pregunta 2}{Enunciado}
\small
Considere:
\begin{itemize}
  \item Modelo de determinación de la renta (aspa keynesiana).
  \item Modelo IS--LM.
  \item Modelo DA--OA.
\end{itemize}

\begin{enumerate}[label=\alph*)]
  \item Compare el impacto de un aumento de $G$ sobre $Y$ en el aspa vs IS--LM.
  \item Compare el impacto de un aumento de $G$ sobre $Y$ en IS--LM vs DA--OA.
  \item Justifique adecuadamente sus respuestas (puede usar gráficos). Excluya casos con curvas perfectamente horizontales o verticales.
\end{enumerate}
\end{frame}

% 2023 P2 Blanco
\begin{frame}{Certamen 2 2023 -- Pregunta 2}
  \vfill
  \centering
  \vfill
\end{frame}

% 2023 P2 Pauta
\begin{frame}{Certamen 2 2023 -- Pregunta 2}{Pauta conceptual}
\small
\textbf{(a) Aspa vs IS--LM}
\begin{itemize}
  \item Aspa: tasa de interés fija, sin mercado de activos.
  \[
  \Delta Y = \frac{1}{1 - c(1-t)}\Delta G.
  \]
  \item IS--LM: $\uparrow G \Rightarrow$ IS derecha, pero $\uparrow i \Rightarrow \downarrow I$ (crowding out).
\end{itemize}
\[
|\Delta Y|_{\text{Aspa}} > |\Delta Y|_{\text{IS--LM}}.
\]

\textbf{(b) IS--LM vs DA--OA}
\begin{itemize}
  \item IS--LM: precios fijos, freno solo por tasas de interés.
  \item DA--OA: $\uparrow G \Rightarrow$ DA derecha, suben $Y$ y $P$; $\uparrow P$ reduce $M/P$ y desplaza LM hacia arriba.
\end{itemize}
\[
|\Delta Y|_{\text{IS--LM}} > |\Delta Y|_{\text{DA--OA}}.
\]
\end{frame}

\begin{frame}{Certamen 2 2023 -- Pregunta 2}{Pauta conceptual (síntesis)}
\small
\textbf{Orden de mayor a menor impacto de $\Delta G$ sobre $Y$}:
\[
|\Delta Y|_{\text{Aspa}} > |\Delta Y|_{\text{IS--LM}} > |\Delta Y|_{\text{DA--OA}}.
\]

\begin{itemize}
  \item Aspa: no hay mercados de activos ni ajuste de precios.
  \item IS--LM: se introduce el crowding out vía tasa de interés.
  \item DA--OA: se agrega además el ajuste de precios y saldos reales.
\end{itemize}
\end{frame}

\end{document}
